% Abstract WAIMLAp 2026 — DINOv2 + LoRA on EuroSAT
% Compile with: tectonic Abstract_DINOv2_EuroSAT_LoRA.tex
\documentclass[10pt,a4paper]{article}
\usepackage[utf8]{inputenc}
\usepackage[T1]{fontenc}
\usepackage{mathptmx}
\usepackage[scaled=0.92]{helvet}
\usepackage{courier}
\usepackage{amssymb}
\usepackage[a4paper,top=2.0cm,bottom=2.0cm,left=1.8cm,right=1.8cm,headsep=0.5cm,footskip=0.8cm]{geometry}
\usepackage{microtype}
\usepackage{booktabs}
\usepackage{tabularx}
\usepackage{caption}
\usepackage{graphicx}
\usepackage{fancyhdr}
\usepackage{xcolor}
\usepackage[hidelinks]{hyperref}
\usepackage{titlesec}

\definecolor{waimlap}{HTML}{4B2A82}

% ---- encabezado y pie del evento
\pagestyle{fancy}
\fancyhf{}
\renewcommand{\headrulewidth}{0pt}
\fancyhead[C]{\sffamily\small\color{waimlap}\textbf{Workshop on Artificial Intelligence and Machine Learning Applications}}
\fancyfoot[C]{\sffamily\footnotesize\color{waimlap}WAIMLAp 2026 \quad|\quad November 18--20, 2026 \quad|\quad Lima, Peru \quad|\quad waimlap.org}

% ---- secciones compactas
\titleformat{\section}{\normalfont\bfseries\small\uppercase}{\thesection.}{0.4em}{}
\titlespacing*{\section}{0pt}{0.75ex}{0.35ex}
\newcommand{\para}[1]{\textbf{#1}\enspace}
\setlength{\parindent}{0pt}
\setlength{\parskip}{0.4ex}
\captionsetup{font=small,labelfont=bf,skip=2.5pt}

\begin{document}

\begin{center}
{\large\bfseries Dissecting and Adapting DINOv2 for Land Cover Classification:\\[1pt]
An Architectural Teardown, Parameter-Efficient Tuning, and Spatial Generalization Study on Sentinel-2 Imagery\par}
\vspace{0.45em}
{\normalsize Alfredo Jiménez, Jhonatan Davila, Italo Ramos\par}
\vspace{0.15em}
{\small Departamento de Ingeniería, Pontificia Universidad Católica del Perú (PUCP), Lima, Peru\\
Grupo de Investigación en Inteligencia Artificial (IA-PUCP)\\
\texttt{alfredo.jf.2004@gmail.com}, \texttt{jcdavilam@pucp.edu.pe}, \texttt{italom.ramosd@pucp.edu.pe}\par}
\end{center}
\vspace{-0.4em}

\section{Introduction}
Self-supervised vision foundation models such as DINOv2~[1], a Vision Transformer (ViT)~[2] pre-trained on 142 million curated natural images, produce general-purpose features that transfer well with a frozen backbone. Satellite imagery, however, is a distinct domain: Sentinel-2 tiles~[3] are acquired from above at 10\,m resolution, dominated by repetitive spatial textures, devoid of canonical foreground objects, and separated across classes by subtle spectral patterns. Linear probing is computationally cheap and preserves pre-trained features but cannot adapt to overhead textural semantics; full fine-tuning updates all 22\,M parameters, risking overfitting and high compute cost. We present a systematic study comprising: (i) an architectural teardown of DINOv2 ViT-S/14; (ii) parameter-efficient Low-Rank Adaptation (LoRA)~[4] on the deepest self-attention blocks; (iii) spatial out-of-distribution evaluation across disjoint geographic regions across Europe (EuroSAT~[5]); and (iv) multi-seed paired statistical verification ($s \in \{0, 1, 42\}$) and architectural ablations.

\section{Methodology}
\para{Dataset and protocols.} We use EuroSAT RGB~[5] (10 land cover classes, 27,000 tiles of $64\times64$ px). Tiles are resized to $224\times224$ px via bicubic interpolation without cropping (avoiding loss of 23\,\% marginal context) and normalized with ImageNet statistics. Following DINOv2 linear evaluation~[1], representations concatenate the [CLS] token and the mean of patch tokens (768 dimensions), yielding a classifier of 7,690 parameters.
To rigorously test spatial transfer, we employ two evaluation protocols:
\begin{itemize}\setlength{\itemsep}{0pt}\setlength{\parskip}{0pt}\setlength{\leftskip}{-1.2em}
\item \textbf{Random stratified split (internal)}: Standard 70/15/15 partition (18,900 train, 4,050 val, 4,050 test), assessing in-distribution performance under identical geographical tile sampling.
\item \textbf{Geographic blocked split (spatial)}: Using UTM coordinates extracted from georeferenced Sentinel-2 multispectral tiles, we perform spatial $K$-means clustering to partition Europe into disjoint geographic regions. Tiles from held-out regions form \texttt{test\_fuera} (unseen geographical domains), while tiles from training clusters form \texttt{test\_dentro}.
\end{itemize}

\para{Phase A: architectural teardown.} DINOv2 ViT-S/14 without registers~[6] comprises 22,056,576 parameters (88.2\,MB in fp32) across 12 Transformer blocks ($d=384$, 6 heads, MLP 1,536). Each block holds 1.78\,M parameters (MLP 67\,\%, attention 33\,\%). Forward compute matches analytical derivation at 12.25\,GFLOPs/img ($224$\,px); quadratic attention constitutes 9.9\,\% of FLOPs, scaling to 30.6\,\% at 448\,px (63.4\,GFLOPs). On a Tesla T4 GPU, latency is 8.89\,ms (batch 1, fp32) and 1.17\,ms/img (batch 128, fp16). Effective receptive field~[7] analysis reveals global gradient coverage from block 1 (50\,\% radius: 69.4\,px in block 1; 85.0\,px in block 12). Token PCA confirms that while early blocks (e.g., block 4) encode spatial grid coordinates, block 12 features align directly with high-level land cover semantics.

\para{Phase B \& C: LoRA intervention and spatial evaluation.} Motivated by semantic concentration in deep layers, blocks 1--10 remain frozen, adding LoRA to query ($Q$) and value ($V$) projections in blocks 11--12. Weight updates take the form $W' = W + \frac{\alpha}{r} BA$, with frozen $W$, $A \sim \mathcal{N}(0, \sigma^2)$, $B=0$ ($r=\alpha=8$), injecting 24,576 parameters (0.11\,\% of backbone; 32,266 including classifier). Training uses AdamW~[8] ($lr=10^{-3}$, $wd=10^{-4}$), batch 64, 8 epochs. In Phase C, we execute 20 controlled runs across 3 random seeds ($s \in \{0, 1, 42\}$) and perform systematic ablations: rank $r \in \{4, 8, 16\}$, depth (2 vs.\ 4 adapted blocks), and projection subsets ($Q$ vs.\ $Q+V$).

\section{Results and Discussion}
\para{Quantitative performance and multi-seed robustness.} Table~1 summarizes test performance across seeds and partition protocols. On the internal random split, LoRA consistently improves macro-F1 from $95.66 \pm 0.16\,\%$ to $97.33 \pm 0.17\,\%$ (+1.67\,pp, error reduction $-32\,\%$). In pairwise exact McNemar tests~[9], discordances heavily favor LoRA across all seeds (seed 42: $b=51, c=97, p=1.9\times10^{-4}$; seed 0: $b=54, c=104, p=8.5\times10^{-5}$; seed 1: $b=24, c=103, p=7.5\times10^{-13}$), confirming statistical significance beyond sampling noise.

\begin{table}[t]
\centering\small
\begin{tabularx}{\linewidth}{@{}lXcccc@{}}
\toprule
Partition & Method & Trainable params & In-Domain Macro-F1 & Out-of-Domain Macro-F1 & Spatial Drop ($\Delta$) \\
\midrule
Internal & Linear probing & 7,690 (0.035\%) & $95.66 \pm 0.16\%$ & --- & --- \\
Internal & LoRA ($r=8$, b11--12) & 32,266 (0.146\%) & $\mathbf{97.33 \pm 0.17\%}$ & --- & --- \\
\midrule
Geographic & Linear probing & 7,690 (0.035\%) & $96.40 \pm 0.18\%$ & $89.25 \pm 0.72\%$ & $-7.15 \pm 0.55$\,pp \\
Geographic & LoRA ($r=8$, b11--12) & 32,266 (0.146\%) & $\mathbf{97.54 \pm 0.17\%}$ & $\mathbf{91.51 \pm 0.38\%}$ & $\mathbf{-6.02 \pm 0.55}$\,pp \\
\bottomrule
\end{tabularx}
\caption{Mean $\pm$ std over 3 seeds ($s \in \{0, 1, 42\}$) on EuroSAT test splits. Spatial drop measures transfer loss from \texttt{test\_dentro} to \texttt{test\_fuera}.}
\end{table}

\para{Spatial generalization and hypothesis testing.} When evaluated on geographically disjoint regions (\texttt{test\_fuera}), both models experience domain shift. Crucially, linear probing drops by $7.15 \pm 0.55$\,pp (to $89.25\,\%$), whereas LoRA drops by only $6.02 \pm 0.55$\,pp (to $91.51\,\%$). Rather than narrowing due to regional overfitting (the pessimistic hypothesis H2), LoRA \emph{widens} its advantage in unseen regions from +1.14\,pp in-domain to \textbf{+2.27\,pp out-of-domain}. Exact McNemar tests on \texttt{test\_fuera} reject the hypothesis of equal accuracy with extreme confidence ($p = 4.2\times10^{-10}$ for seed 42; $p = 2.9\times10^{-11}$ for seed 0; $p = 6.1\times10^{-7}$ for seed 1; discordances $c/b > 2.2$). Thus, LoRA acts as a domain-adaptive spatial regularizer.

\para{Class breakdown and ablations.} Linear structures show marked recovery under LoRA: Highway rises by +5.85\,pp and River by +3.71\,pp, eliminating over 90\,\% of River$\rightarrow$Highway confusions (Fig.~1). Residual errors reside largely among spectrally ambiguous vegetation classes (e.g., Pasture vs.\ AnnualCrop at 64\,px without NIR/SWIR).
Ablation analysis ($s=42$) reveals three architectural insights:
\begin{enumerate}\setlength{\itemsep}{0pt}\setlength{\parskip}{0pt}\setlength{\leftskip}{-1.0em}
\item \textbf{Rank invariance}: $r=4$ achieves $97.27\,\%$ internal F1 and $90.54\,\%$ spatial F1, statistically matching $r=8$ ($97.22\,\%$ and $90.60\,\%$) while $r=16$ yields no additional gain ($96.86\,\%$ and $90.97\,\%$), demonstrating low intrinsic adaptation rank.
\item \textbf{Depth benefit}: Expanding LoRA to 4 blocks (blocks 8--11) achieves the highest global performance: \textbf{97.97\,\%} internal F1 (98.05\,\% acc) and \textbf{92.27\,\%} spatial F1 (94.88\,\% acc), providing superior resilience against spatial domain shift.
\item \textbf{Projection efficiency}: Adapting only $Q$ achieves $97.24\,\%$ internal F1 and $90.60\,\%$ spatial F1, matching $Q+V$ while halving LoRA parameters to 12,288 (0.055\% of backbone).
\end{enumerate}

\begin{figure}[b]
\centering
\includegraphics[width=0.88\linewidth]{fig_predicciones.png}
\caption{Representative test predictions with softmax confidence. LoRA corrects thin linear features (River vs.\ Highway) while maintaining calibrated predictions.}
\end{figure}

\section{Conclusions}
Adapting the final self-attention blocks of DINOv2 with low-rank matrices (0.11\,\% of parameters) bridges the domain gap between natural and orbital imagery without inference overhead upon weight merging. Rigorous geographic splitting across 3 seeds demonstrates that LoRA significantly mitigates spatial performance decay, outperforming linear probing by +2.27\,pp in unseen geographical domains ($p < 10^{-6}$). Ablation confirms that rank $r=4$ suffices and that adapting 4 deep blocks provides maximal spatial transfer resilience.

\section*{References}
\vspace{-0.25em}
{\footnotesize
\begin{enumerate}\setlength{\itemsep}{0pt}\setlength{\parskip}{0pt}\setlength{\leftskip}{-0.6em}
\renewcommand{\labelenumi}{[\arabic{enumi}]}
\item M. Oquab \emph{et al.}, ``DINOv2: Learning robust visual features without supervision,'' \emph{TMLR}, 2024.
\item A. Dosovitskiy \emph{et al.}, ``An image is worth 16x16 words: Transformers for image recognition at scale,'' in \emph{Proc. ICLR}, 2021.
\item M. Drusch \emph{et al.}, ``Sentinel-2: ESA's optical mission for GMES,'' \emph{Remote Sens. Environ.}, vol.~120, pp.~25--36, 2012.
\item E. J. Hu \emph{et al.}, ``LoRA: Low-rank adaptation of large language models,'' in \emph{Proc. ICLR}, 2022.
\item P. Helber \emph{et al.}, ``EuroSAT: Land use and land cover classification,'' \emph{IEEE JSTARS}, vol.~12, no.~7, 2019.
\item T. Darcet \emph{et al.}, ``Vision transformers need registers,'' in \emph{Proc. ICLR}, 2024.
\item W. Luo \emph{et al.}, ``Understanding the effective receptive field in deep CNNs,'' in \emph{Proc. NeurIPS}, 2016.
\item I. Loshchilov and F. Hutter, ``Decoupled weight decay regularization,'' in \emph{Proc. ICLR}, 2019.
\item Q. McNemar, ``Note on the sampling error of the difference between correlated proportions,'' \emph{Psychometrika}, vol.~12, 1947.
\item D. R. Roberts \emph{et al.}, ``Cross-validation strategies for spatial or temporal structure,'' \emph{Ecography}, vol.~40, 2017.
\end{enumerate}}

\end{document}